\section{Exactly summable commutator relations}\label{sec:special}
The universal series is infinite, but algebraic relations can remove almost all of it. The following proofs establish global exponential identities, not just small-neighborhood Taylor expansions. For Banach algebras all exponentials exist; for finite-dimensional Lie groups the same proofs use one-parameter subgroups and their logarithmic derivatives. Complex scalar parameters are understood in a complex Banach algebra or complex Lie group; in a real Lie group the parameters and linear combinations must be real.

\subsection{Commuting elements and central commutators}
If $[X,Y]=0$, the binomial proof in \cref{lem:exp-log} gives
\begin{equation}\label{eq:commuting}
 e^Xe^Y=e^{X+Y}
\end{equation}
without a smallness restriction.

\begin{theorem}[Central-commutator or disentangling identity]\label{thm:central}
Suppose $C=[X,Y]$ commutes with both $X$ and $Y$. Then, for every scalar parameter $t$ for which the exponentials are interpreted,
\begin{align}
 e^{tX}Ye^{-tX}&=Y+tC,\label{eq:central-conjugation}\\
 e^{tX}e^{tY}&=\exp\!\left(t(X+Y)+\tfrac12t^2C\right).\label{eq:central-t}
\end{align}
In particular,
\begin{equation}\label{eq:central-bch}
 e^Xe^Y=e^{X+Y+C/2},\qquad
 e^{X+Y}=e^Xe^Ye^{-C/2},\qquad
 e^{X/2}e^Ye^{X/2}=e^{X+Y}.
\end{equation}
\end{theorem}
\begin{proof}
All terms of \cref{eq:campbell} after the first commutator vanish, proving \cref{eq:central-conjugation}. If $g(t)=e^{tX}e^{tY}$, then
\begin{equation}\label{eq:central-ode}
 g'(t)=\bigl(X+e^{tX}Ye^{-tX}\bigr)g(t)
       =(X+Y+tC)g(t),\qquad g(0)=I.
\end{equation}
The exponent $K(t)=t(X+Y)+t^2C/2$ commutes with its derivative, so differentiating its exponential gives the same equation and initial condition. Uniqueness proves \cref{eq:central-t}. Taking $t=1$ gives the first identity in \cref{eq:central-bch}, and commuting the central exponential gives the second. Applying the first identity twice with $X/2,Y,X/2$ cancels the two central corrections, proving the symmetric identity.
\end{proof}
Centrality is needed only relative to the Lie subalgebra generated by $X,Y$, not relative to a larger ambient algebra.

A useful complete finite-product extension is the following. If all commutators $[X_i,X_j]$ commute with every $X_k$, then
\begin{equation}\label{eq:central-many}
 \prod_{j=1}^{m}e^{X_j}
 =\exp\!\left(\sum_{j=1}^{m}X_j+
       \frac12\sum_{1\leq i<j\leq m}[X_i,X_j]\right).
\end{equation}
The proof is induction on $m$ using \cref{thm:central}; the new commutator is the sum of the pairwise commutators with the last element. This is the all-factors form behind the Weyl and displacement multiplication laws.

\subsection{An eigenvector of the adjoint action}
\begin{theorem}[{The relation $[X,Y]=sY$}]\label{thm:eigenline}
Suppose $[X,Y]=sY$ with $s\in\CC$. Define
\[
 q_s(t)=\begin{cases}(1-e^{-st})/s,&s\ne0,\\t,&s=0.\end{cases}
\]
Then for every scalar $c$,
\begin{equation}\label{eq:eigen-factor}
 e^{t(X+cY)}=e^{tX}e^{c q_s(t)Y}.
\end{equation}
Consequently, when $s\notin2\pi\ii\mathbb Z\setminus\{0\}$,
\begin{equation}\label{eq:eigen-bch}
 \boxed{\displaystyle e^Xe^Y=\exp\!\left(X+\frac{s}{1-e^{-s}}Y\right),}
\end{equation}
with the coefficient interpreted as $1$ at $s=0$. For every $s$, including resonant values,
\begin{equation}\label{eq:eigen-braid}
 e^Xe^Y=e^{e^sY}e^X,\qquad e^Xe^Ye^{-X}=e^{e^sY}.
\end{equation}
\end{theorem}
\begin{proof}
Equation \cref{eq:campbell} gives $e^{tX}Ye^{-tX}=e^{st}Y$. The right logarithmic derivative of
$F(t)=e^{tX}e^{c q_s(t)Y}$ is therefore
\[
 F'F^{-1}=X+c q_s'(t)e^{st}Y=X+cY,
\]
because $q_s'(t)=e^{-st}$. This proves \cref{eq:eigen-factor} by uniqueness of the differential equation with $F(0)=I$. At $t=1$, choose $c=1/q_s(1)=s/(1-e^{-s})$ when this is defined. The apparent singularity at zero is removable, and the nonzero zeros of $1-e^{-s}$ are exactly the excluded values. The braiding formulas follow directly from conjugating the exponential and require no division by $q_s(1)$.
\end{proof}
Every bracket containing at least two $Y$'s vanishes under this relation. To see this, evaluate any bracket tree. A branch containing one or more $Y$'s, when nonzero, is a multiple of $Y$. At the first node combining two such branches the bracket is zero. A branch consisting only of $X$'s is zero unless it has length one. Thus only the terms linear in $Y$ survive, and formally
\begin{equation}\label{eq:eigen-series}
 Z=X+\sum_{n\geq0}b_n^+s^nY.
\end{equation}
This also follows from \cref{eq:H1}, but the differential-equation proof establishes \cref{eq:eigen-bch} even when that Taylor series diverges.

If $Y\ne0$, \cref{eq:bernoulli-zeta} shows that \cref{eq:eigen-series} converges for $|s|<2\pi$ and diverges for $|s|\geq2\pi$: its even-index terms have magnitudes proportional to
$2\zeta(2m)(|s|/(2\pi))^{2m}$ and therefore do not tend to zero in the latter range. The closed expression \cref{eq:eigen-bch} can still be valid there.

For an explicit illustration take
\begin{equation}\label{eq:eigen-matrices}
 X=\begin{pmatrix}s&0\\0&0\end{pmatrix},\qquad
 Y=\begin{pmatrix}0&1\\0&0\end{pmatrix}.
\end{equation}
Direct exponentiation gives
\[
 e^Xe^Y=\begin{pmatrix}e^s&e^s\\0&1\end{pmatrix},\qquad
 e^{X+cY}=\begin{pmatrix}e^s&c(e^s-1)/s\\0&1\end{pmatrix},
\]
with the usual continuous value at $s=0$. If $s=3\pi$, the choice $c=h(s)$ works although the BCH series diverges. If $s=2\pi\ii k\ne0$, every $e^{X+cY}$ is $I$, whereas $e^Xe^Y=I+Y$. This shows what fails at resonance: not the braiding identity, but the possibility of a logarithm of the particular form $X+cY$. In this example a different logarithm, namely $Y$, still exists.
